% Epilogue. Blake's voice at the other seam. Target ~900 words.
%
% Seed paragraph (from the existing web closing) at the bottom of this file.
% Expand it: the Tassafaronga moment should land here with its proper weight,
% as a moment the reader lives through, not a fact that is reported. The
% torpedo took 183 of his shipmates and most of Joan's replies down with
% the ship. Their absence in this book has a date.
%
% After Tassafaronga, Gene came home. He and Joan married. They were
% married for forty-nine years.
%
% Style guidance: present at the seam, not woven throughout. Let the
% epilogue be the place where the gravity of the project lands.

\cleardoublepage
\thispagestyle{laempty}
\vspace*{0.10\textheight}
\begin{center}
    \eyebrow{Epilogue}\par
    \vspace{0.6em}
    \brassrule[36pt]\par
\end{center}
\vspace{2em}

\noindent
% =====================================================================
% TODO[blake]: REPLACE THIS WITH YOUR EXPANDED EPILOGUE.
% =====================================================================
%
% Below is the seed paragraph from the website's closing. It works as
% a one-paragraph epilogue, but the gravity it carries is what the body
% of this book has earned the right to. Expand it to ~900 words. Live
% through Tassafaronga; do not summarize it. Live the homecoming. Then
% say what forty-nine years of marriage was, in whatever spare form
% suits the voice.
% ---------------------------------------------------------------------

\noindent
\textsc{[\,seed paragraph for expansion. Replace.\,]}
\vspace{0.4em}

\noindent
\itshape
Gene's letters to Joan continued through the war. Less than a year
after the last letter of 1940, on the morning of December 7, 1941,
he was at Pearl Harbor. A year after that, off Tassafaronga in the
Solomon Islands, a Japanese torpedo struck the New Orleans and tore
away one hundred and fifty feet of her bow. One hundred and eighty-three
of his shipmates went down with it, along with most of Joan's letters
back. Gene came home in 1943. He and Joan were married for forty-nine
years.
\upshape

\vspace{1.4em}
\brassornament

\vspace{1em}
\noindent
\textsc{[\,placeholder \,---\, two to four paragraphs in Blake's voice after the
italic line above. Where the silence in the archive comes from. What
Joan kept. What got lost. What Blake found, and what he didn't find.
The closing should be one line.\,]}
\vfill

\begin{flushright}
    {\handfont\Large B.M.}\par
\end{flushright}
